\documentclass[11pt]{article}
\usepackage[a4paper,margin=2.1cm]{geometry}
\usepackage[T1]{fontenc}
\usepackage{lmodern}
\usepackage{microtype}
\usepackage{amsmath,amssymb,mathtools,bm}
\usepackage{graphicx,booktabs,xcolor,enumitem}
\usepackage[most]{tcolorbox}
\usepackage{caption,float}
\captionsetup{font=small,labelformat=empty,skip=2pt}
\definecolor{acc}{HTML}{C0392B}
\definecolor{muted}{HTML}{666666}
\setlength{\parindent}{0pt}
\setlength{\parskip}{5pt}
\setlist{nosep,leftmargin=1.2em}
\newcommand{\vb}[1]{\boldsymbol{#1}}
\newcommand{\eps}{\varepsilon}
\newcommand{\br}[1]{\langle #1 |}
\newcommand{\kt}[1]{| #1 \rangle}
\newtcolorbox{plain}[1]{colback=acc!4,colframe=acc!60,boxrule=0.5pt,arc=2pt,left=6pt,right=6pt,top=4pt,bottom=4pt,title={#1},coltitle=black,colbacktitle=acc!12,fonttitle=\bfseries}
\pagestyle{plain}

\begin{document}

{\Large\bfseries From a volume to a boundary: the paper's mode expansion for a 2-D particle}\par
\vspace{2pt}
{\color{muted} Companion note to Chen, Bergman and Sivan (arXiv 1711.00335; Phys.\ Rev.\ Applied 11, 044018, 2019), equations (21) and (26).
Prepared for Claudio Silvestre Castro, 4 November 2026. Equation numbers in brackets like (9) are the paper's; primed ones like (21$'$) are this note's.}

\section*{1\quad The paper in three lines}

The paper writes the field as its equation (11), $\vb E = \vb E_0 + (\eps_i-\eps_b)\,\hat\Gamma\hat\theta\,\vb E$, where $\hat\Gamma\hat\theta$ means ``integrate the background Green function against the field, over the particle only''. Moving the unknown to one side:
\begin{equation}
M(\eps_i)\,\vb E=\vb E_0,\qquad M(\eps_i)=1-(\eps_i-\eps_b)\,\hat\Gamma\hat\theta .
\tag{A}
\end{equation}
The particle permittivity $\eps_i$ appears once, as a number in front of a fixed operator: \textbf{$M$ is a straight line in $\eps_i$.}

A mode is a value $\eps_m$ and a field $\vb E_m$ with $M(\eps_m)\vb E_m=0$, that is $\hat\Gamma\hat\theta\,\vb E_m=s_m\vb E_m$ with $s_m=1/(\eps_m-\eps_b)$, equations (9)--(10). Each mode is scaled so that $\br{E_m}\hat\theta\kt{E_m}=\int_{\text{particle}}\vb E_m\cdot\vb E_m\,d\vb r=1$, equation (16), with no complex conjugate, equation (34). The two results are
\begin{align*}
\vb E &= \vb E_0+\sum_m \vb E_m\,\frac{\eps_i-\eps_b}{\eps_m-\eps_i}\,\br{E_m}\hat\theta\kt{E_0}, \tag{21}\\
\bar{\bar G}(\vb r,\vb r') &= \bar{\bar G}_0(\vb r,\vb r')+\frac1{k^2}\sum_m \frac{\eps_i-\eps_b}{(\eps_m-\eps_i)(\eps_m-\eps_b)}\,\vb E_m(\vb r)\otimes\vb E_m(\vb r'). \tag{26}
\end{align*}

\section*{2\quad The same problem written on the boundary}

For a long cylinder, Maxwell's equations reduce to one field component $\psi$ along the axis ($E_y$ or $H_y$). Instead of the field inside the particle, the unknowns are two functions on the boundary curve: the field $\varphi=\psi$ and its outward slope $\chi=\partial\psi/\partial n$. With $n$ boundary points they form one vector $x=[\varphi;\chi]$ of $2n$ numbers (Fig.~1).

\begin{figure}[H]
\centering
\includegraphics[width=\linewidth]{fig1.pdf}
\caption*{\textbf{Fig.\ 1.} (a) The unknowns sit on the boundary. (b) The matrix: the outside rows use the background Green function $g_b$ (any slab included) and never see $\eps_i$; only the inside rows change with $\eps_i$.}
\end{figure}
\vspace{-6pt}

Green's identity gives two sets of $n$ equations each (Fig.~1b):
\begin{itemize}
\item \textbf{outside rows}, built with $g_b$: \quad $A\varphi+B\chi=b$, where $b_j=\psi_0(\text{point } j)$ is the incident field on the boundary;
\item \textbf{inside rows}, built with the outgoing wave $H_0^{(1)}(k_i|\vb r-\vb r'|)$ at the inside wavenumber $k_i=k\sqrt{\eps_i}$: \quad $C(\eps_i)\varphi+\eta\,D(\eps_i)\chi=0$, with $\eta=\eps_i/\eps_b$ for $H_y$ and $\eta=1$ for $E_y$.
\end{itemize}
\begin{equation}
M(\eps_i)\,x=\begin{bmatrix}b\\0\end{bmatrix},\qquad
\psi(\vb r)=\psi_0(\vb r)+P(\vb r)\,x\quad\text{(field at any outside point $\vb r$)}.
\tag{B}
\end{equation}
$P(\vb r)$ is a fixed row of numbers: the field at $\vb r$ radiated by the boundary values (pysie2d's \texttt{eval\_field}).

\textbf{The one difference that matters:} $\eps_i$ now sits \emph{inside} the Green function, through $\sqrt{\eps_i}$. So $M$ is a curve in $\eps_i$, not a straight line.

\section*{3\quad The weight of one mode, in four lines}

A mode is a value $\eps_m$ and a vector $v_m$ with $M(\eps_m)v_m=0$: boundary values of a field that exists with no incident wave. Each mode also has a row partner $w_m$ with $w_m^{T}M(\eps_m)=0$. In the paper the partner is the mode itself ($\br{E_m}\hat\theta$, equation 34); on the boundary $M$ is not symmetric, so $w_m$ is a different vector, computed alongside $v_m$.

Multiply (B) on the left by $w_m^T$. Because $w_m^TM(\eps_m)=0$, we can subtract it for free:
\[
w_m^T\big[M(\eps_i)-M(\eps_m)\big]\,x=w_m^T b .
\]
Close to $\eps_m$: $M(\eps_i)-M(\eps_m)\approx(\eps_i-\eps_m)\,M'(\eps_m)$, with $M'=dM/d\eps_i$, and one mode dominates, $x\approx c_m v_m$. Therefore
\begin{equation}
c_m(\eps_i)=\frac{w_m^T b}{(\eps_i-\eps_m)\,N_m},\qquad N_m=w_m^T M'(\eps_m)\,v_m .
\tag{C}
\end{equation}
$N_m$ replaces the paper's scaling (16). Only the inside rows of $M'$ are non-zero.

\textbf{Check against the paper.} For (A): $M'=-\hat\Gamma\hat\theta$ and $w_m=\br{E_m}\hat\theta$, so $N_m=-s_m\br{E_m}\hat\theta\kt{E_m}=-s_m$. Then (C) gives
$c_m=\br{E_m}\hat\theta\kt{E_0}\,(\eps_m-\eps_b)/(\eps_m-\eps_i)$: exactly the weight in the paper's equation (19). For a straight line both ``$\approx$'' are exact, because the partners of different modes do not mix ($w_m^TM'v_n=0$ for $n\neq m$).

\section*{4\quad What is left over}

\begin{plain}{Keldysh's theorem, in plain words}
Let $M(\eps)$ be a matrix whose entries change smoothly with complex $\eps$. Near each mode value $\eps_m$ (a mode that is not repeated), its inverse looks like a single fraction:
\[
M(\eps)^{-1}\approx \frac{v_m\,w_m^T}{(\eps-\eps_m)\,N_m}.
\]
Subtract these fractions for all modes in a region of the complex $\eps$ plane. What remains has no blow-ups in that region and changes smoothly with $\eps$. {\color{muted}(M.\ V.\ Keldysh, 1951.)}
\end{plain}

The same idea with ordinary numbers, $1/f(\eps)$ instead of a matrix inverse:
\begin{itemize}
\item a straight line, $f=a(\eps-\eps_1)$: $\;1/f=1/[a(\eps-\eps_1)]$ exactly, so nothing is left over;
\item a curve, $f=e^{\eps}-1$, which is zero at $\eps=2\pi i\ell$ for every whole number $\ell$:
\[
\frac1{e^{\eps}-1}=-\frac12+\sum_{\ell}\frac1{\eps-2\pi i \ell}\qquad(\text{adding $\ell$ and $-\ell$ together}).
\]
\end{itemize}
Even with every zero included, $-\tfrac12$ is left over.

So, for any $M$:
\begin{equation}
x(\eps_i)=\sum_m c_m(\eps_i)\,v_m+x_{\text{smooth}}(\eps_i).
\tag{D}
\end{equation}

\textbf{In the paper (straight line) the leftover does not depend on $\eps_i$.} From (19)--(20), the leftover is $\vb E_0-\sum_m \vb E_m\br{E_m}\hat\theta\kt{E_0}$, with no $\eps_i$ in it; it vanishes inside the particle (equation 15). Equation (21) folds this constant back into $\vb E_0$, using
\[
\frac{\eps_m-\eps_b}{\eps_m-\eps_i}=1+\frac{\eps_i-\eps_b}{\eps_m-\eps_i}.
\]
That is where (21)'s factor $(\eps_i-\eps_b)/(\eps_m-\eps_i)$ comes from.

\textbf{On the boundary (curve) the leftover changes with $\eps_i$}, smoothly. It also collects every mode you did not include. A few exact solves pin it down.

\section*{5\quad The boundary versions of (21) and (26)}

Put (C) and (D) into (B). With $F_m(\vb r)=P(\vb r)\,v_m$, the field of mode $m$ outside the particle:
\begin{equation}
\psi(\vb r)=\psi_0(\vb r)+\sum_m F_m(\vb r)\,\frac{w_m^T b}{(\eps_i-\eps_m)\,N_m}+P(\vb r)\,x_{\text{smooth}}(\eps_i).
\tag{21$'$}
\end{equation}
For a line source at $\vb r'$, the incident field is $g_b(\vb r,\vb r')$ and $b_j=g_b(\text{point } j,\vb r')$. The overlap $w_m^Tb$ becomes $\tilde F_m(\vb r')=\sum_j w_{m,j}\,g_b(\text{point } j,\vb r')$:
\begin{equation}
G(\vb r,\vb r')=g_b(\vb r,\vb r')+\sum_m \frac{F_m(\vb r)\,\tilde F_m(\vb r')}{(\eps_i-\eps_m)\,N_m}+P(\vb r)\,x_{\text{smooth}}(\eps_i;\vb r').
\tag{26$'$}
\end{equation}
$\tilde F_m$ plays the part of the paper's ``adjoint'' mode $\vb E^\dagger_m$. Equation (34) says $\vb E^\dagger_m=\vb E_m$. Here, the response at $\vb r$ to a source at $\vb r'$ equals the response at $\vb r'$ to a source at $\vb r$; that forces $\tilde F_m=\kappa_m F_m$ for a constant $\kappa_m$. Checked with pysie2d for two modes of a particle with no symmetry, at five points: equal to 1 part in $10^{14}$. So (26$'$) has the paper's shape, with the term $F_m(\vb r)F_m(\vb r')\,\kappa_m/[(\eps_i-\eps_m)N_m]$.

\begin{center}
\small
\renewcommand{\arraystretch}{1.25}
\begin{tabular}{@{}p{3.1cm}p{6.1cm}p{6.6cm}@{}}
\toprule
 & \textbf{Volume (paper)} & \textbf{Boundary (2-D particle, pysie2d)}\\
\midrule
unknown & $\vb E(\vb r)$ inside the particle & $\varphi,\chi$ at $n$ boundary points\\
how $\eps_i$ enters & number $(\eps_i-\eps_b)$ in front of $\hat\Gamma\hat\theta$ & inside the Green function, via $k\sqrt{\eps_i}$\\
mode & $\hat\Gamma\hat\theta\vb E_m=s_m\vb E_m$ & $M(\eps_m)v_m=0$\\
partner & $\br{E_m}\hat\theta$, the mode itself & $w_m$ with $w_m^TM(\eps_m)=0$, a different vector\\
scaling & $\br{E_m}\hat\theta\kt{E_m}=1$ & $N_m=w_m^TM'(\eps_m)v_m$\\
overlap with source & $\int_{\text{particle}}\vb E_m\cdot\vb E_0\,d\vb r$ & $\sum_j w_{m,j}\,\psi_0(\text{point } j)$, outside rows only\\
weight & $(\eps_i-\eps_b)/(\eps_m-\eps_i)$ & $1/[(\eps_i-\eps_m)N_m]$\\
leftover & constant in $\eps_i$, folded into $\vb E_0$ & smooth in $\eps_i$; fitted from a few exact solves\\
layered background & layered Green function inside every mode calculation (9) & $g_b$ only in the outside rows, at the real wavelength, built once\\
\bottomrule
\end{tabular}
\end{center}

\section*{6\quad How well it works}

\begin{figure}[H]
\centering
\includegraphics[width=\linewidth]{fig2.pdf}
\caption*{\textbf{Fig.\ 2.} Rounded triangle 140\,nm above a 400\,nm slab of $\varepsilon = 4$ in silica, 1550\,nm, $H_y$. Amplitude of the slab's guided wave under a plane wave, versus the particle permittivity.}
\end{figure}
\vspace{-6pt}

The six mode terms alone (dashed) miss by 19--39\,\%. Modelling the leftover as $a+b\,\eps_i+c\,\eps_i^2$, from three exact solves in each of the three ranges (red; the middle range is shaded), brings the error to 0.14--1.5\,\%, except near $\eps_i=-\eps_b$ (28\,\%). There, many sharp modes crowd together and no short formula captures them; the paper meets the same crowding.

\end{document}
